\section{Threats to Validity}\label{sec:threats}

\emph{Internal.} On CIC-IoT2023 the same author identified the attackers, wrote the labeling rule and made all design choices. The labels were revised once after error analysis (the sixth DDoS source), towards the model's predictions; a later blind audit in both directions found one more error, whose correction also favors the model (Benign FPR 1.3\% instead of 3.8\%; Section~\ref{sec:labels}). The audit's rater is a program, not a second person, and four of its rules were fixed after its first run (one a post-hoc criterion); we release its blind 975-flow sample for a manual check. Our CIC-IoT2023 results are scored against our own labels; IoT-23 supports the cost of file-name labels with the authors' labels but does not validate ours, and its file-name labels depend on our family mapping.

\emph{Construct.} Macro-F1 weighs tiny classes (Mirai, MITM-ARP) like large ones and, on IoT-23, missing classes; we report per-class results with host-pair counts. Without MITM-ARP, which flow features cannot see by the label's definition, macro-F1 is 0.826 instead of 0.725. False-block rates depend on the traffic mix and on the counting unit (60\,s host-pair windows live, host pairs over 5.9\,h in the replay), and the 60\,s reset hides some false blocks (Section~\ref{sec:response}).

\emph{External.} Labels and splits were audited on two datasets, but the live and response checks only on CIC-IoT2023, with one capture session per attack class, so they measure generalization within a session, not across networks. We do not test across datasets, where results are known to be asymmetric~\cite{ali2026siamxbert}. We audit one model: a random forest and gradient boosting score similarly on the temporal split (Section~\ref{sec:setup}), but we did not repeat the label and response checks with them. Large CIC-IoT2023 captures were trimmed, so a temporal split covers only their start. Live results hold only at the achieved replay rates, on one laptop under WSL2, and for floods that did not fill the capture buffer. Replay is open-loop, so ``packets dropped'' is counterfactual: drops cause no retransmissions, and attackers do not adapt.

\emph{Statistical conclusion.} The enforced settings ran once, and each pinned CPU budget three times. The SSH-to-RTSP change rests on three host pairs. Beyond it, the CIC-IoT2023 gap ($+0.03$ to $+0.04$) excludes zero by host pair and attacker host but not by attacker device or capture, units with little power. Per-class intervals over at most six host pairs are indicative only.

\section{Conclusion}\label{sec:conclusion}

Each check exposed a problem that the aggregate score of its stage could not show, from labels to blocking, and once in our own live results, where flow splitting left macro-F1 unchanged. We recommend that evaluations of flow-based NIDS (i) report each attack capture's background share and a leave-one-capture-out FPR; (ii) decompose a random--temporal gap per class and host pair, with class presence and the boundary-spanning share, before calling it drift; (iii) match live to offline flows one by one on every run, since an offline run flushes every flow at the end of the file; and (iv) score blocking in false blocks per host-pair window, with persistent counters and verdicts timed at flow export. Releasing CIC-IoT2023 flows with addresses and timestamps would let users relabel them. Future work includes a blocking criterion that weighs attack volume, not only flow verdicts, longer live runs, a second live dataset, early classification~\cite{pekar2025early} and edge hardware.

\subsection{Artifact Availability}

The repository \url{https://github.com/muhammadusamahoyrr/EdgeFlow-IDS} (MIT license) holds the code, attacker evidence, configuration, result files and audit scripts. The scripts recompute the offline numbers from the datasets; the live numbers also need the live flow logs, which are not in the repository. All random choices (split, trees, NSGA-II, bootstrap) use seed 0; live runs are timing-dependent. The NFStream fix is included (\nolinkurl{src/patch_nfstream.py}, diff in \nolinkurl{patches/}). The version used for this paper is the tag \texttt{v1.0-paper}.
% PLACEHOLDER for the author (do not compile in until minted): add the Zenodo DOI of tag
% v1.0-paper, e.g. "..., archived at Zenodo (DOI <to be minted>)".
The CIC-IoT2023 and IoT-23 data are available from their authors~\cite{neto2023ciciot2023,garcia2020iot23}.
